\section{十五世纪的冲突、灾害与司法证据}

战争、灾荒、疾病、赈恤、案狱和边地暴力并不是彼此独立的材料领域。军报和判牍往往由胜者、官府或审讯机关形成；灾情、医疗与慈善记录也常以请赈、捐输、工程或死亡数字的形式出现。它们能显示国家和地方组织如何命名风险与行动，却不能把受害者、病者、被征发者或被定罪者的经验全部还原。

\subsection{军政冲突与动员}

\eventanchor{ML-1445-002}{明·1445（正统十年）}
\paragraph{1445（正统十年）：浙江干旱鼠疫肆虐} 此一账本锚点将冲突、风险或法律处置置于具体时段。它所能可靠支持的判断是编年候选的年序可由官修与后出材料交叉；具体月日、参与者、战果或数字必须回查所列材料，不能将单一叙事当作定论。材料链由，因而同时保存了命令、报送、传记、地方记忆或跨政权叙事的不同立场。问题形成的条件是自然风险与地方报告促成朝廷或地方的应对记录；随后产生的可见影响是赈济、迁徙和损失的实际范围仍须以受灾地区材料分别核验。对于医疗救助、慈善网络、普通人的伤亡与案狱处境，材料往往尤为稀薄。译写候选以编年材料定位；实录记载反映朝廷选择，崇祯期须另以奏疏、地方及敌对政权材料交叉。

\eventanchor{ML-1446-001}{明·1446（正统十一年）一月}
\paragraph{1446（正统十一年）一月：陆川平勉之战：司仁法被处决} 此一账本锚点将冲突、风险或法律处置置于具体时段。它所能可靠支持的判断是编年候选的年序可由官修与后出材料交叉；具体月日、参与者、战果或数字必须回查所列材料，不能将单一叙事当作定论。材料链由，因而同时保存了命令、报送、传记、地方记忆或跨政权叙事的不同立场。问题形成的条件是前期边防、地方武装或跨区域资源调动的压力推动了该行动；随后产生的可见影响是它改变了后续调兵、补给或地方秩序的条件；战果和伤亡须分开核对。对于医疗救助、慈善网络、普通人的伤亡与案狱处境，材料往往尤为稀薄。译写候选以编年材料定位；实录记载反映朝廷选择，崇祯期须另以奏疏、地方及敌对政权材料交叉。

\eventanchor{ML-1446-002}{明·1446（正统十一年）}
\paragraph{1446（正统十一年）：江南洪水泛滥} 此一账本锚点将冲突、风险或法律处置置于具体时段。它所能可靠支持的判断是编年候选的年序可由官修与后出材料交叉；具体月日、参与者、战果或数字必须回查所列材料，不能将单一叙事当作定论。材料链由，因而同时保存了命令、报送、传记、地方记忆或跨政权叙事的不同立场。问题形成的条件是自然风险与地方报告促成朝廷或地方的应对记录；随后产生的可见影响是赈济、迁徙和损失的实际范围仍须以受灾地区材料分别核验。对于医疗救助、慈善网络、普通人的伤亡与案狱处境，材料往往尤为稀薄。译写候选以编年材料定位；实录记载反映朝廷选择，崇祯期须另以奏疏、地方及敌对政权材料交叉。

\eventanchor{ML-1446-003}{明·1446（正统十一年）}
\paragraph{1446（正统十一年）：学校、科举和礼制的本年政策材料为社会流动和国家知识秩序提供可检索锚点} 此一账本锚点将冲突、风险或法律处置置于具体时段。它所能可靠支持的判断是来源导向的年度桥接条目：本年材料可定位相关政务线索；具体月日、发文机关、地域和执行结果仍须逐项回查，不能以制度文本或奏报替代实际效果。材料链由，因而同时保存了命令、报送、传记、地方记忆或跨政权叙事的不同立场。问题形成的条件是制度、教育或知识传播的需要推动了相关行动或文本形成；随后产生的可见影响是它提供制度和传播史的锚点，不能据此推断社会整体接受。对于医疗救助、慈善网络、普通人的伤亡与案狱处境，材料往往尤为稀薄。桥接条目只标示可回查的年度问题与材料链；实录有中央选择性，崇祯期尤其需要奏疏、地方、朝鲜及满文材料交叉。

\eventanchor{ML-1447-001}{明·1447（正统十二年）}
\paragraph{1447（正统十二年）：叶宗柳与浙江银矿工人叛乱} 此一账本锚点将冲突、风险或法律处置置于具体时段。它所能可靠支持的判断是编年候选的年序可由官修与后出材料交叉；具体月日、参与者、战果或数字必须回查所列材料，不能将单一叙事当作定论。材料链由，因而同时保存了命令、报送、传记、地方记忆或跨政权叙事的不同立场。问题形成的条件是财政征解、交通工程或货币支付的压力使相关措施进入政务；随后产生的可见影响是其法定安排不等于地方执行，后续须以地域材料追踪负担与成效。对于医疗救助、慈善网络、普通人的伤亡与案狱处境，材料往往尤为稀薄。译写候选以编年材料定位；实录记载反映朝廷选择，崇祯期须另以奏疏、地方及敌对政权材料交叉。

\eventanchor{ML-1447-002}{明·1447（正统十二年）}
\paragraph{1447（正统十二年）：苏北饥荒袭来} 此一账本锚点将冲突、风险或法律处置置于具体时段。它所能可靠支持的判断是编年候选的年序可由官修与后出材料交叉；具体月日、参与者、战果或数字必须回查所列材料，不能将单一叙事当作定论。材料链由，因而同时保存了命令、报送、传记、地方记忆或跨政权叙事的不同立场。问题形成的条件是自然风险与地方报告促成朝廷或地方的应对记录；随后产生的可见影响是赈济、迁徙和损失的实际范围仍须以受灾地区材料分别核验。对于医疗救助、慈善网络、普通人的伤亡与案狱处境，材料往往尤为稀薄。译写候选以编年材料定位；实录记载反映朝廷选择，崇祯期须另以奏疏、地方及敌对政权材料交叉。

\subsection{灾害、医疗与赈恤}

\eventanchor{ML-1447-003}{明·1447（正统十二年）}
\paragraph{1447（正统十二年）：灾情上报、赈恤或河工材料标记本年地方风险，损失与救济效果需回到受灾地区核实} 此一账本锚点将冲突、风险或法律处置置于具体时段。它所能可靠支持的判断是来源导向的年度桥接条目：本年材料可定位相关政务线索；具体月日、发文机关、地域和执行结果仍须逐项回查，不能以制度文本或奏报替代实际效果。材料链由，因而同时保存了命令、报送、传记、地方记忆或跨政权叙事的不同立场。问题形成的条件是自然风险与地方报告促成朝廷或地方的应对记录；随后产生的可见影响是赈济、迁徙和损失的实际范围仍须以受灾地区材料分别核验。对于医疗救助、慈善网络、普通人的伤亡与案狱处境，材料往往尤为稀薄。桥接条目只标示可回查的年度问题与材料链；实录有中央选择性，崇祯期尤其需要奏疏、地方、朝鲜及满文材料交叉。

\eventanchor{ML-1448-001}{明·1448（正统十三年）三月}
\paragraph{1448（正统十三年）三月：邓茂奇与一群佃农在闽赣边境西北叛乱} 此一账本锚点将冲突、风险或法律处置置于具体时段。它所能可靠支持的判断是编年候选的年序可由官修与后出材料交叉；具体月日、参与者、战果或数字必须回查所列材料，不能将单一叙事当作定论。材料链由，因而同时保存了命令、报送、传记、地方记忆或跨政权叙事的不同立场。问题形成的条件是继承、官僚协调或皇权运作的压力使问题进入朝廷程序；随后产生的可见影响是它影响后续人事与政策议程；动机和派系标签不能由单一叙事裁定。对于医疗救助、慈善网络、普通人的伤亡与案狱处境，材料往往尤为稀薄。译写候选以编年材料定位；实录记载反映朝廷选择，崇祯期须另以奏疏、地方及敌对政权材料交叉。

\eventanchor{ML-1448-002}{明·1448（正统十三年）十二月}
\paragraph{1448（正统十三年）十二月：明军杀了叶宗六，但他的叛军毫发无伤，继续向南撤退，围困滁州} 此一账本锚点将冲突、风险或法律处置置于具体时段。它所能可靠支持的判断是编年候选的年序可由官修与后出材料交叉；具体月日、参与者、战果或数字必须回查所列材料，不能将单一叙事当作定论。材料链由，因而同时保存了命令、报送、传记、地方记忆或跨政权叙事的不同立场。问题形成的条件是前期边防、地方武装或跨区域资源调动的压力推动了该行动；随后产生的可见影响是它改变了后续调兵、补给或地方秩序的条件；战果和伤亡须分开核对。对于医疗救助、慈善网络、普通人的伤亡与案狱处境，材料往往尤为稀薄。译写候选以编年材料定位；实录记载反映朝廷选择，崇祯期须另以奏疏、地方及敌对政权材料交叉。

\eventanchor{ML-1448-003}{明·1448（正统十三年）}
\paragraph{1448（正统十三年）：1448年黄河洪水：黄河堤决口} 此一账本锚点将冲突、风险或法律处置置于具体时段。它所能可靠支持的判断是编年候选的年序可由官修与后出材料交叉；具体月日、参与者、战果或数字必须回查所列材料，不能将单一叙事当作定论。材料链由，因而同时保存了命令、报送、传记、地方记忆或跨政权叙事的不同立场。问题形成的条件是自然风险与地方报告促成朝廷或地方的应对记录；随后产生的可见影响是赈济、迁徙和损失的实际范围仍须以受灾地区材料分别核验。对于医疗救助、慈善网络、普通人的伤亡与案狱处境，材料往往尤为稀薄。译写候选以编年材料定位；实录记载反映朝廷选择，崇祯期须另以奏疏、地方及敌对政权材料交叉。

\eventanchor{ML-1448-004}{明·1448（正统十三年）}
\paragraph{1448（正统十三年）：干旱和蝗灾袭击中国西北地区} 此一账本锚点将冲突、风险或法律处置置于具体时段。它所能可靠支持的判断是编年候选的年序可由官修与后出材料交叉；具体月日、参与者、战果或数字必须回查所列材料，不能将单一叙事当作定论。材料链由，因而同时保存了命令、报送、传记、地方记忆或跨政权叙事的不同立场。问题形成的条件是自然风险与地方报告促成朝廷或地方的应对记录；随后产生的可见影响是赈济、迁徙和损失的实际范围仍须以受灾地区材料分别核验。对于医疗救助、慈善网络、普通人的伤亡与案狱处境，材料往往尤为稀薄。译写候选以编年材料定位；实录记载反映朝廷选择，崇祯期须另以奏疏、地方及敌对政权材料交叉。

\eventanchor{ML-1448-005}{明·1448（正统十三年）}
\paragraph{1448（正统十三年）：江南大旱} 此一账本锚点将冲突、风险或法律处置置于具体时段。它所能可靠支持的判断是编年候选的年序可由官修与后出材料交叉；具体月日、参与者、战果或数字必须回查所列材料，不能将单一叙事当作定论。材料链由，因而同时保存了命令、报送、传记、地方记忆或跨政权叙事的不同立场。问题形成的条件是自然风险与地方报告促成朝廷或地方的应对记录；随后产生的可见影响是赈济、迁徙和损失的实际范围仍须以受灾地区材料分别核验。对于医疗救助、慈善网络、普通人的伤亡与案狱处境，材料往往尤为稀薄。译写候选以编年材料定位；实录记载反映朝廷选择，崇祯期须另以奏疏、地方及敌对政权材料交叉。

\subsection{审案、处分与地方程序}

\eventanchor{ML-1449-001}{明·1449（正统十四年）三月}
\paragraph{1449（正统十四年）三月：陆川平绵之战：明军因藏匿司继发而入侵蒙阳，司继发又逃脱} 此一账本锚点将冲突、风险或法律处置置于具体时段。它所能可靠支持的判断是编年候选的年序可由官修与后出材料交叉；具体月日、参与者、战果或数字必须回查所列材料，不能将单一叙事当作定论。材料链由，因而同时保存了命令、报送、传记、地方记忆或跨政权叙事的不同立场。问题形成的条件是前期边防、地方武装或跨区域资源调动的压力推动了该行动；随后产生的可见影响是它改变了后续调兵、补给或地方秩序的条件；战果和伤亡须分开核对。对于医疗救助、慈善网络、普通人的伤亡与案狱处境，材料往往尤为稀薄。译写候选以编年材料定位；实录记载反映朝廷选择，崇祯期须另以奏疏、地方及敌对政权材料交叉。

\eventanchor{ML-1449-002}{明·1449（正统十四年）五月}
\paragraph{1449（正统十四年）五月：邓茂奇叛军被击败} 此一账本锚点将冲突、风险或法律处置置于具体时段。它所能可靠支持的判断是编年候选的年序可由官修与后出材料交叉；具体月日、参与者、战果或数字必须回查所列材料，不能将单一叙事当作定论。材料链由，因而同时保存了命令、报送、传记、地方记忆或跨政权叙事的不同立场。问题形成的条件是前期边防、地方武装或跨区域资源调动的压力推动了该行动；随后产生的可见影响是它改变了后续调兵、补给或地方秩序的条件；战果和伤亡须分开核对。对于医疗救助、慈善网络、普通人的伤亡与案狱处境，材料往往尤为稀薄。译写候选以编年材料定位；实录记载反映朝廷选择，崇祯期须另以奏疏、地方及敌对政权材料交叉。

\eventanchor{ML-1449-003}{明·1449（正统十四年）七月}
\paragraph{1449（正统十四年）七月：图木危机：北元事实上的统治者瓦剌也先太师入侵明朝} 此一账本锚点将冲突、风险或法律处置置于具体时段。它所能可靠支持的判断是编年候选的年序可由官修与后出材料交叉；具体月日、参与者、战果或数字必须回查所列材料，不能将单一叙事当作定论。材料链由，因而同时保存了命令、报送、传记、地方记忆或跨政权叙事的不同立场。问题形成的条件是继承、官僚协调或皇权运作的压力使问题进入朝廷程序；随后产生的可见影响是它影响后续人事与政策议程；动机和派系标签不能由单一叙事裁定。对于医疗救助、慈善网络、普通人的伤亡与案狱处境，材料往往尤为稀薄。译写候选以编年材料定位；实录记载反映朝廷选择，崇祯期须另以奏疏、地方及敌对政权材料交叉。

\eventanchor{ML-1449-004}{明·1449（正统十四年）八月4日}
\paragraph{1449（正统十四年）八月4日：土木危机：正统皇帝出京亲自迎战也先太史} 此一账本锚点将冲突、风险或法律处置置于具体时段。它所能可靠支持的判断是编年候选的年序可由官修与后出材料交叉；具体月日、参与者、战果或数字必须回查所列材料，不能将单一叙事当作定论。材料链由，因而同时保存了命令、报送、传记、地方记忆或跨政权叙事的不同立场。问题形成的条件是继承、官僚协调或皇权运作的压力使问题进入朝廷程序；随后产生的可见影响是它影响后续人事与政策议程；动机和派系标签不能由单一叙事裁定。对于医疗救助、慈善网络、普通人的伤亡与案狱处境，材料往往尤为稀薄。译写候选以编年材料定位；实录记载反映朝廷选择，崇祯期须另以奏疏、地方及敌对政权材料交叉。

\eventanchor{ML-1449-005}{明·1449（正统十四年）八月30日}
\paragraph{1449（正统十四年）八月30日：土木危机：明军后卫被击败} 此一账本锚点将冲突、风险或法律处置置于具体时段。它所能可靠支持的判断是编年候选的年序可由官修与后出材料交叉；具体月日、参与者、战果或数字必须回查所列材料，不能将单一叙事当作定论。材料链由，因而同时保存了命令、报送、传记、地方记忆或跨政权叙事的不同立场。问题形成的条件是前期边防、地方武装或跨区域资源调动的压力推动了该行动；随后产生的可见影响是它改变了后续调兵、补给或地方秩序的条件；战果和伤亡须分开核对。对于医疗救助、慈善网络、普通人的伤亡与案狱处境，材料往往尤为稀薄。译写候选以编年材料定位；实录记载反映朝廷选择，崇祯期须另以奏疏、地方及敌对政权材料交叉。

\eventanchor{ML-1449-006}{明·1449（正统十四年）八月}
\paragraph{1449（正统十四年）八月：叶宗流的叛军被击败} 此一账本锚点将冲突、风险或法律处置置于具体时段。它所能可靠支持的判断是编年候选的年序可由官修与后出材料交叉；具体月日、参与者、战果或数字必须回查所列材料，不能将单一叙事当作定论。材料链由，因而同时保存了命令、报送、传记、地方记忆或跨政权叙事的不同立场。问题形成的条件是前期边防、地方武装或跨区域资源调动的压力推动了该行动；随后产生的可见影响是它改变了后续调兵、补给或地方秩序的条件；战果和伤亡须分开核对。对于医疗救助、慈善网络、普通人的伤亡与案狱处境，材料往往尤为稀薄。译写候选以编年材料定位；实录记载反映朝廷选择，崇祯期须另以奏疏、地方及敌对政权材料交叉。

\subsection{边地暴力与行政后果}

\eventanchor{ML-1449-007}{明·1449（正统十四年）九月1日}
\paragraph{1449（正统十四年）九月1日：土木危机：明军全军覆没，正统皇帝被也先太师俘虏} 此一账本锚点将冲突、风险或法律处置置于具体时段。它所能可靠支持的判断是编年候选的年序可由官修与后出材料交叉；具体月日、参与者、战果或数字必须回查所列材料，不能将单一叙事当作定论。材料链由，因而同时保存了命令、报送、传记、地方记忆或跨政权叙事的不同立场。问题形成的条件是前期边防、地方武装或跨区域资源调动的压力推动了该行动；随后产生的可见影响是它改变了后续调兵、补给或地方秩序的条件；战果和伤亡须分开核对。对于医疗救助、慈善网络、普通人的伤亡与案狱处境，材料往往尤为稀薄。译写候选以编年材料定位；实录记载反映朝廷选择，崇祯期须另以奏疏、地方及敌对政权材料交叉。

\eventanchor{ML-1449-008}{明·1449（正统十四年）九月23日}
\paragraph{1449（正统十四年）九月23日：朱祁钰即位景泰皇帝} 此一账本锚点将冲突、风险或法律处置置于具体时段。它所能可靠支持的判断是编年候选的年序可由官修与后出材料交叉；具体月日、参与者、战果或数字必须回查所列材料，不能将单一叙事当作定论。材料链由，因而同时保存了命令、报送、传记、地方记忆或跨政权叙事的不同立场。问题形成的条件是继承、官僚协调或皇权运作的压力使问题进入朝廷程序；随后产生的可见影响是它影响后续人事与政策议程；动机和派系标签不能由单一叙事裁定。对于医疗救助、慈善网络、普通人的伤亡与案狱处境，材料往往尤为稀薄。译写候选以编年材料定位；实录记载反映朝廷选择，崇祯期须另以奏疏、地方及敌对政权材料交叉。

\eventanchor{ML-1449-009}{明·1449（正统十四年）十月27日}
\paragraph{1449（正统十四年）十月27日：也先太史围攻北京未能攻下，五天后撤军} 此一账本锚点将冲突、风险或法律处置置于具体时段。它所能可靠支持的判断是编年候选的年序可由官修与后出材料交叉；具体月日、参与者、战果或数字必须回查所列材料，不能将单一叙事当作定论。材料链由，因而同时保存了命令、报送、传记、地方记忆或跨政权叙事的不同立场。问题形成的条件是前期边防、地方武装或跨区域资源调动的压力推动了该行动；随后产生的可见影响是它改变了后续调兵、补给或地方秩序的条件；战果和伤亡须分开核对。对于医疗救助、慈善网络、普通人的伤亡与案狱处境，材料往往尤为稀薄。译写候选以编年材料定位；实录记载反映朝廷选择，崇祯期须另以奏疏、地方及敌对政权材料交叉。

\eventanchor{ML-1449-010}{明·1449（正统十四年）}
\paragraph{1449（正统十四年）：黄河堤坝再次决堤导致河道略有改道} 此一账本锚点将冲突、风险或法律处置置于具体时段。它所能可靠支持的判断是编年候选的年序可由官修与后出材料交叉；具体月日、参与者、战果或数字必须回查所列材料，不能将单一叙事当作定论。材料链由，因而同时保存了命令、报送、传记、地方记忆或跨政权叙事的不同立场。问题形成的条件是财政征解、交通工程或货币支付的压力使相关措施进入政务；随后产生的可见影响是其法定安排不等于地方执行，后续须以地域材料追踪负担与成效。对于医疗救助、慈善网络、普通人的伤亡与案狱处境，材料往往尤为稀薄。译写候选以编年材料定位；实录记载反映朝廷选择，崇祯期须另以奏疏、地方及敌对政权材料交叉。

\eventanchor{ML-1450-001}{明·1450（景泰元年）九月19日}
\paragraph{1450（景泰元年）九月19日：正统皇帝被释放回到北京，被景泰皇帝软禁} 此一账本锚点将冲突、风险或法律处置置于具体时段。它所能可靠支持的判断是编年候选的年序可由官修与后出材料交叉；具体月日、参与者、战果或数字必须回查所列材料，不能将单一叙事当作定论。材料链由，因而同时保存了命令、报送、传记、地方记忆或跨政权叙事的不同立场。问题形成的条件是继承、官僚协调或皇权运作的压力使问题进入朝廷程序；随后产生的可见影响是它影响后续人事与政策议程；动机和派系标签不能由单一叙事裁定。对于医疗救助、慈善网络、普通人的伤亡与案狱处境，材料往往尤为稀薄。译写候选以编年材料定位；实录记载反映朝廷选择，崇祯期须另以奏疏、地方及敌对政权材料交叉。

\eventanchor{ML-1450-002}{明·1450（景泰元年）}
\paragraph{1450（景泰元年）：贵州、湖广瑶族、苗族起义} 此一账本锚点将冲突、风险或法律处置置于具体时段。它所能可靠支持的判断是编年候选的年序可由官修与后出材料交叉；具体月日、参与者、战果或数字必须回查所列材料，不能将单一叙事当作定论。材料链由，因而同时保存了命令、报送、传记、地方记忆或跨政权叙事的不同立场。问题形成的条件是继承、官僚协调或皇权运作的压力使问题进入朝廷程序；随后产生的可见影响是它影响后续人事与政策议程；动机和派系标签不能由单一叙事裁定。对于医疗救助、慈善网络、普通人的伤亡与案狱处境，材料往往尤为稀薄。译写候选以编年材料定位；实录记载反映朝廷选择，崇祯期须另以奏疏、地方及敌对政权材料交叉。
